%% tfg.tex — Documento principal. Aquí solo se decide QUÉ entra y en QUÉ orden.
%% Para quitar un apartado, pon % al principio de su línea.
\documentclass{estilo/memoria}           % formato de la memoria (no tocar)
%% datos.tex — Datos de tu trabajo. Escribe entre las llaves { }.
%% Los campos marcados como (opcional) pueden quedar vacíos, p. ej. \cotutor{},
%% o borrarse sin que nada falle. Los datos de la universidad (nombre, logo,
%% titulación, márgenes…) están en estilo/institucion.tex.
%% Desde Estudio TFG también puedes editarlos en el panel «Datos del trabajo».

\titulo{Título de tu Trabajo Fin de Grado}
\autor{Nombre Apellido1 Apellido2}
\email{nombre.apellido@alu.uclm.es}             % (opcional)
\tutor{Dr. Nombre Apellido}
\cotutor{}                                      % (opcional)
\departamento{Tecnologías y Sistemas de Información}  % (opcional) sale como «Departamento de …»
\intensificacion{Ingeniería del Software}      % (opcional) tu tecnología específica
\fecha{junio}{2026}                             % (opcional) mes y año de entrega; si falta, la fecha actual
\ciudad{Ciudad Real}                            % (opcional) sale en la página de créditos
\palabrasClave{palabra clave 1, palabra clave 2, palabra clave 3}  % (opcional) para el Resumen
\keywords{keyword 1, keyword 2, keyword 3}      % (opcional) las mismas, en inglés, para el Abstract

%% Ajustes del documento
\idioma{espanol}              % espanol | ingles (títulos automáticos: «Capítulo», «Índice»…)
\formato{impresion}           % impresion (doble cara, enlaces negros) | pantalla (una cara, enlaces en color)
\modo{borrador}               % borrador (marca «BORRADOR» al pie y notas \todo en rojo) | final (un \todo da error)
\estiloBibliografia{ieee}     % ieee (numérico, orden de cita) | apa (autor-año) | numeric | authoryear | alphabetic
\licencia{reservados}         % reservados | cc-by | cc-by-sa | cc-by-nc-sa | ninguna (texto de la página de créditos)
\atribucion{si}               % si | no (página final que cita la clase original de ARCO)
                            % tus datos: título, autor, tutor, idioma…

\begin{document}

%% ---- Parte inicial (páginas en números romanos) ----
\frontmatter                             % empieza la parte inicial
\portada                                 % portada exterior
\portadainterior                         % portada interior con tutores
\paginaderechos                          % créditos: autor, contacto y licencia
\tribunal                                % hoja para el tribunal
\input{0-inicio/resumen}                 % resumen en español
\input{0-inicio/abstract}                % resumen en inglés
\input{0-inicio/agradecimientos}         % agradecimientos (opcional)
\indices                                 % índice general y, si hay, de figuras, tablas y listados
\input{0-inicio/acronimos}               % lista de acrónimos

%% ---- Capítulos (números arábigos) ----
\mainmatter                              % empieza el cuerpo del trabajo
\chapter{Introducción}\label{cap:introduccion}
% GUÍA — Qué va en este capítulo:
%  - El contexto del trabajo y el problema concreto que resuelve.
%  - Por qué merece la pena resolverlo (motivación, a quién beneficia).
%  - Una idea general de la solución propuesta, sin entrar en detalle.
%  - Al final, un párrafo que describa cómo se organiza el resto de la memoria.

Sustituye este párrafo por la presentación de tu \ac{TFG}: el contexto en el
que surge, el problema que aborda y por qué es relevante.

La memoria se organiza así: el capítulo~\ref{cap:objetivos} fija los objetivos,
el capítulo~\ref{cap:antecedentes} repasa los antecedentes, el
capítulo~\ref{cap:metodologia} describe la metodología, el
capítulo~\ref{cap:resultados} presenta los resultados y el
capítulo~\ref{cap:conclusiones} recoge las conclusiones.
    % capítulo 1: introducción
\chapter{Objetivos}\label{cap:objetivos}
% GUÍA — Qué va en este capítulo:
%  - El objetivo general del trabajo, en una o dos frases claras.
%  - Los objetivos específicos (parciales) que lo concretan, mejor en una lista
%    numerada y redactados en infinitivo: «Diseñar…», «Evaluar…».
%  - Cada objetivo debe poder comprobarse al final (en Resultados/Conclusiones).
%  - Si aplica: requisitos, restricciones o limitaciones de partida.

\section{Objetivo general}

Sustituye este párrafo por el objetivo general de tu trabajo.

\section{Objetivos específicos}

\begin{enumerate}
  \item Primer objetivo específico.
  \item Segundo objetivo específico.
\end{enumerate}
       % capítulo 2: objetivos
\chapter{Antecedentes}\label{cap:antecedentes}
% GUÍA — Qué va en este capítulo (también llamado «estado del arte»):
%  - Los conceptos y tecnologías necesarios para entender el trabajo.
%  - Qué soluciones o trabajos previos existen y en qué se quedan cortos.
%  - Toda afirmación que no sea tuya debe ir con su referencia (\cite{…}).
%  - Termina mostrando el hueco que tu trabajo viene a cubrir.

Sustituye este párrafo por la revisión de los trabajos y tecnologías
relacionados con tu propuesta, citando siempre sus fuentes.
    % capítulo 3: antecedentes / estado del arte
\chapter{Metodología}\label{cap:metodologia}
% GUÍA — Qué va en este capítulo:
%  - La metodología de trabajo seguida (p. ej. un proceso iterativo o ágil)
%    y cómo se ha aplicado en la práctica: iteraciones, planificación, roles.
%  - Las herramientas, lenguajes y recursos hardware/software utilizados.
%  - Si procede, el diseño de la solución (arquitectura, modelos, diagramas).
%  - Este capítulo incluye ejemplos de figura, tabla, cita, listado y acrónimo
%    para que copies su estructura; bórralos cuando ya no los necesites.

Sustituye este párrafo por la descripción de cómo has llevado a cabo el
trabajo.

\section{Ejemplos para copiar}

% --- CITA: \cite{clave}, donde «clave» es la primera palabra de la entrada
%     en bibliografia.bib. El ~ evita que la cita quede sola en otra línea.
El ciclo de trabajo se inspira en los métodos iterativos descritos por
Sommerville~\cite{sommerville2016}.

% --- ACRÓNIMO: \ac{SIGLA}; la sigla debe estar definida en 0-inicio/acronimos.tex.
Los modelos del diseño se expresan en \ac{UML}.

% --- REFERENCIAS CRUZADAS: \ref{etiqueta} da el número de una figura, tabla,
%     listado o capítulo que tenga \label{etiqueta}.
La figura~\ref{fig:ciclo} resume el proceso, la tabla~\ref{tab:herramientas}
recoge las herramientas empleadas y el listado~\ref{lst:hola} muestra un
fragmento de código.

% --- FIGURA: el archivo va en la carpeta figuras/ (PDF, PNG o JPG) y se
%     escribe sin la carpeta. width=… fija el ancho (aquí, 80 % del texto).
\begin{figure}[htbp]
  \centering
  \includegraphics[width=0.8\textwidth]{ejemplo}
  \caption{Ciclo de trabajo iterativo seguido en el proyecto.}
  \label{fig:ciclo}
\end{figure}

% --- TABLA: columnas l (izquierda), c (centro), r (derecha); & separa
%     celdas y \\ termina la fila. \toprule, \midrule y \bottomrule son las
%     líneas horizontales.
\begin{table}[htbp]
  \centering
  \caption{Herramientas utilizadas.}
  \label{tab:herramientas}
  \begin{tabular}{lll}
    \toprule
    \textbf{Herramienta} & \textbf{Uso}           & \textbf{Versión} \\
    \midrule
    Python               & Desarrollo             & 3.12 \\
    Git                  & Control de versiones   & 2.45 \\
    \LaTeX{}             & Redacción de la memoria & TeX Live 2026 \\
    \bottomrule
  \end{tabular}
\end{table}

% --- LISTADO DE CÓDIGO: language=… activa el resaltado (Python, Java, C, C++,
%     SQL, HTML…). El código se copia tal cual entre begin y end.
\begin{lstlisting}[language=Python, caption={Ejemplo de función en Python.}, label={lst:hola}]
def saludar(nombre):
    """Devuelve un saludo para la persona indicada."""
    return f"Hola, {nombre}"
\end{lstlisting}
     % capítulo 4: metodología
\chapter{Resultados}\label{cap:resultados}
% GUÍA — Qué va en este capítulo:
%  - Lo que se ha obtenido al aplicar la metodología: el sistema construido,
%    su evolución por iteraciones o los experimentos realizados.
%  - Las pruebas o la evaluación, con datos (tablas, gráficas) y su análisis.
%  - Relaciona cada resultado con los objetivos del capítulo~\ref{cap:objetivos}.

Sustituye este párrafo por los resultados obtenidos y su análisis.
      % capítulo 5: resultados
\chapter{Conclusiones}\label{cap:conclusiones}
% GUÍA — Qué va en este capítulo:
%  - Repasa cada objetivo y explica en qué medida se ha cumplido.
%  - La justificación de las competencias de tu tecnología específica que
%    has puesto en práctica en el trabajo.
%  - Las líneas de trabajo futuro: mejoras y ampliaciones posibles.
%  - Si quieres, una valoración personal breve.

Sustituye este párrafo por las conclusiones de tu trabajo.

\section{Trabajo futuro}

Describe aquí las mejoras y ampliaciones que quedan abiertas.
    % capítulo 6: conclusiones

%% ---- Parte final ----
\bibliography{bibliografia}              % referencias (de bibliografia.bib)
\appendix                                % a partir de aquí, anexos (A, B…)
\include{2-anexos/a-anexo}               % anexo A

\end{document}
